\documentclass[11pt]{article}
\usepackage[a4paper,margin=1in]{geometry}
\usepackage[T1]{fontenc}
\usepackage{lmodern}
\usepackage{amsmath,amssymb,amsthm,booktabs,enumitem,mathtools,microtype}
\usepackage[numbers,sort&compress]{natbib}
\usepackage[hidelinks]{hyperref}

\newtheorem{theorem}{Theorem}[section]
\newtheorem{proposition}[theorem]{Proposition}
\newtheorem{lemma}[theorem]{Lemma}
\newtheorem{corollary}[theorem]{Corollary}
\theoremstyle{definition}
\newtheorem{definition}[theorem]{Definition}
\theoremstyle{remark}
\newtheorem{remark}[theorem]{Remark}

\newcommand{\T}{\{-1,0,+1\}}
\newcommand{\mob}{\mu}
\newcommand{\liou}{\lambda}
\newcommand{\dd}{\delta}
\newcommand{\thh}{\vartheta}
\newcommand{\kap}{\kappa}
\newcommand{\Glog}{G_{\log}}
\newcommand{\Zq}{\mathbb Z/q\mathbb Z}
\newcommand{\ind}{\mathbf 1}
\DeclareMathOperator{\lcm}{lcm}

\hypersetup{
  pdftitle={Fixed-Lag Terminal-Log Möbius Diagonalization and the Square-Divisor Capacity Landscape},
  pdfauthor={RH research program},
  pdfsubject={Fixed-lag terminal logarithmic Mobius diagonalization and exact safe-table capacity},
  pdfkeywords={Mobius function, logarithmic Chowla, squarefree pairs, finite-state capacity}
}

\title{Fixed-Lag Terminal-Log M\"obius Diagonalization and the Square-Divisor Capacity Landscape}
\author{RH research program}
\date{August 10, 2026}

\begin{document}
\maketitle

\begin{abstract}
We study terminal logarithmic averages along fixed periodic phases.  First,
for any two fixed positive-leading affine forms with arbitrary nonzero
determinant, we prove full-M\"obius cancellation with a fixed periodic mask.
The proof inserts finite squarefree cutoffs, combines congruences with an
\(\lcm\) modulus without assuming the cutoff divisors coprime, extracts
contents, and applies Tao's fixed nonparallel Liouville theorem.  A
determinant-free tail estimate completes the limits in the order: cutoff,
terminal limit, then cutoff removal.

This cancellation diagonalizes every fixed finite collection of distinct
M\"obius shifts for polynomials of total degree at most two.  Separately,
for each fixed lag \(h\ge1\), it diagonalizes coordinatewise-biquadratic
polynomials in \((\mob_0(n-h),\mob(n))\).  A local finite-CRT argument gives
the exact two-site density \(\thh^{(h)}_{q,r}\), including collisions modulo
\(p\) and \(p^2\).  Projecting \(512\) truth tables to eight actions and
charging every plus phase to its forced-empty predecessor yields, for all
fixed \(q,h\ge1\),
\[
 \Glog(q,h)=\frac6{\pi^2}-\frac{\kap_h}{2},\qquad
 \kap_h=\prod_{p^2\mid h}(1-p^{-2})
          \prod_{p^2\nmid h}(1-2p^{-2}).
\]
Thus squarefree lags are exactly the maximizers.  The range is
\(3/\pi^2<\Glog(q,h)\le6/\pi^2-\kap_\star/2\), where
\(\kap_\star=\prod_p(1-2p^{-2})\), and its infimum is not attained.  Every
limit is formed for fixed data before the finite table maximum; no growing
period, growing lag, ordinary Ces\`aro, or multi-lag degree-three claim is
made.
\end{abstract}

\noindent\textbf{Keywords:} M\"obius function; terminal logarithmic
average; affine two-point cancellation; squarefree pairs; finite-state
capacity.

\section{Definitions and results}

Write \(\mob_0(m)=\mob(m)\) for \(m\ge1\) and \(\mob_0(m)=0\) for
\(m\le0\).  A terminal clock is a function satisfying
\begin{equation}
 1\le \omega(X)\le X,\qquad \omega(X)\longrightarrow\infty.
 \label{eq:clock}
\end{equation}
All periods, shifts, affine forms, coefficients, and tables below are fixed
before \(X\to\infty\).

For \(q\ge1\) and \(r\in\Zq\), define the one-site and lag-\(h\) two-site
densities
\begin{align}
 \dd_{q,r}&=\lim_{N\to\infty}\frac1N
  \sum_{\substack{n\le N\\n\equiv r\ (q)}}\mob(n)^2,
 \label{eq:delta-def}\\
 \thh^{(h)}_{q,r}&=\lim_{N\to\infty}\frac1N
  \sum_{\substack{n\le N\\n\equiv r\ (q)}}
  \mob_0(n-h)^2\mob(n)^2.
 \label{eq:theta-def}
\end{align}
Their local formulas will be proved in Section~\ref{sec:density}.

\begin{theorem}[Finite-shift total-degree-two diagonalization]
\label{thm:finite-shift}
Fix \(q\), distinct integers \(a_1,\ldots,a_m\), and phase polynomials
\[
 P_r(x_1,\ldots,x_m)=c_\varnothing(r)+
 \sum_i c_i(r)x_i+\sum_i c_{ii}(r)x_i^2+
 \sum_{i<j}c_{ij}(r)x_ix_j
\]
of total degree at most two.  For every clock \eqref{eq:clock},
\begin{align}
 &\frac1{\log\omega(X)}\sum_{X/\omega(X)<n\le X}
 \frac{P_{n\bmod q}(\mob_0(n-a_1),\ldots,\mob_0(n-a_m))}{n}
 \notag\\
 &\hspace{35mm}\longrightarrow
 \sum_{r\bmod q}\left\{\frac{c_\varnothing(r)}q+
 \sum_i c_{ii}(r)\dd_{q,r-a_i}\right\}.
 \label{eq:finite-shift}
\end{align}
Every linear and off-diagonal quadratic channel vanishes.
\end{theorem}

The next statement is deliberately separate: \(x^2z^2\) has total degree
four but coordinatewise degree two.

\begin{theorem}[Fixed-lag coordinatewise-biquadratic compiler]
\label{thm:biquadratic}
Fix integers \(q,h\ge1\) and
\(P_r(x,z)=\sum_{0\le i,j\le2}c_{ij}(r)x^iz^j\).  Then
\begin{align}
 &\frac1{\log\omega(X)}\sum_{X/\omega(X)<n\le X}
 \frac{P_{n\bmod q}(\mob_0(n-h),\mob(n))}{n}
 \notag\\
 &\quad\longrightarrow\sum_{r\bmod q}\left\{
 \frac{c_{00}(r)}q+c_{20}(r)\dd_{q,r-h}
 +c_{02}(r)\dd_{q,r}+c_{22}(r)\thh^{(h)}_{q,r}\right\}.
 \label{eq:biquadratic}
\end{align}
The other five coefficient channels vanish.
\end{theorem}

For a truth-table specialization, let \(f_r:\T^2\to\{-1,+1\}\), put
\(E_r=\{(x,z):f_r(x,z)=+1\}\), and call the phase tuple \emph{lag-\(h\)
safe} if no \((x,z)\in E_r\) and \((z,w)\in E_{r+h}\) coexist.  Denote the
finite set of safe tuples by \(\mathcal A_{q,h}\), and set
\begin{equation}
 L_{q,h}(f;\omega)=\lim_{X\to\infty}\frac1{\log\omega(X)}
 \sum_{X/\omega(X)<n\le X}
 \frac{\mob(n)f_{n\bmod q}(\mob_0(n-h),\mob(n))}{n}.
 \label{eq:score}
\end{equation}
Theorem~\ref{thm:biquadratic} will show that this limit exists and is
independent of \(\omega\).  Only after these fixed-table limits are formed,
define
\begin{equation}
 \Glog(q,h)=\max_{f\in\mathcal A_{q,h}}|L_{q,h}(f;\omega)|.
 \label{eq:capacity-def}
\end{equation}

\begin{theorem}[Capacity and lag spectrum]
\label{thm:capacity}
For every pair of fixed integers \(q,h\ge1\),
\begin{equation}
 \boxed{\Glog(q,h)=\frac6{\pi^2}-\frac{\kap_h}{2}},\qquad
 \kap_h=\prod_{p^2\mid h}(1-p^{-2})
          \prod_{p^2\nmid h}(1-2p^{-2}).
 \label{eq:capacity}
\end{equation}
The value does not depend on \(q\).  Constant table \(36\) attains the
positive sign and its input reflection, table \(72\), attains the negative
sign.  With \(\kap_\star=\prod_p(1-2p^{-2})\),
\begin{equation}
 \frac3{\pi^2}<\Glog(q,h)\le
 \frac6{\pi^2}-\frac{\kap_\star}{2}.
 \label{eq:range}
\end{equation}
Equality on the right holds exactly for squarefree \(h\); the left endpoint
is the unattained infimum over \(h\ge1\).
\end{theorem}

\section{Terminal analytic tools}

\begin{lemma}[Terminal Abel transfer]
\label{lem:abel}
Let \((a_n)\) be bounded and
\(A(t)=\sum_{n\le t}a_n=\alpha t+o(t)\).  Then every clock
\eqref{eq:clock} satisfies
\begin{equation}
 \frac1{\log\omega(X)}\sum_{X/\omega(X)<n\le X}\frac{a_n}{n}
 \longrightarrow\alpha.
 \label{eq:abel}
\end{equation}
\end{lemma}

\begin{proof}
Put \(Y=X/\omega(X)\) and \(B(t)=A(t)-\alpha t\).  Partial summation gives
\[
 \sum_{Y<n\le X}\frac{a_n}{n}
 =\alpha\log(X/Y)+\frac{B(X)}X-\frac{B(Y)}Y
   +\int_Y^X\frac{B(t)}{t^2}\,dt+O(1).
\]
For any \(\varepsilon>0\), choose \(T\) with
\(|B(t)|\le\varepsilon t\) for \(t\ge T\).  If \(Y\ge T\), the normalized
error is at most \(\varepsilon+o(1)\).  If \(Y<T\), the part below \(T\)
is bounded and \(\log X/\log\omega(X)\to1\).  The same bound follows.
Let \(\varepsilon\downarrow0\).
\end{proof}

\begin{theorem}[Terminal full-M\"obius cancellation]
\label{thm:full-mu}
Let
\[
 D(n)=a_1n+b_1,\qquad V(n)=a_2n+b_2,
 \quad a_1,a_2\in\mathbb N,\quad b_1,b_2\in\mathbb Z,
\]
and suppose \(\Delta=a_1b_2-a_2b_1\ne0\).  For every fixed bounded
periodic \(\rho\) and every clock \eqref{eq:clock},
\begin{equation}
 \sum_{X/\omega(X)<n\le X}
 \frac{\mob(D(n))\mob(V(n))\rho(n)}n=o(\log\omega(X)).
 \label{eq:full-mu}
\end{equation}
Values before both forms become positive are omitted, changing the sum by
\(O(1)\).
\end{theorem}

\begin{proof}
Let \(P\) be fixed, larger than every prime divisor of \(a_1a_2\), and
\[
 S_P(m)=\prod_{p\le P}(1-\ind_{p^2\mid m})
 =\sum_{k\mid P\#}\mob(k)\ind_{k^2\mid m}.
\]
Since \(\mob(m)=\liou(m)\ind_{m\ \mathrm{squarefree}}\), first replace the
two M\"obius factors by
\(\liou(D)\liou(V)S_P(D)S_P(V)\).  Let \(M\) be a period of \(\rho\).
Expanding both cutoffs and the fixed periodic mask leaves finitely many systems
\[
 n\equiv s\pmod M,\qquad k^2\mid D(n),\qquad \ell^2\mid V(n).
\]
Each compatible system is a finite disjoint union of residue classes
\(n=r+Lt\) modulo \(L\), where
\begin{equation}
 L=\lcm(M,k^2,\ell^2)\le M(P\#)^2.
 \label{eq:lcm}
\end{equation}
No assumption \((k,\ell)=1\) is used.

On such a class, set
\[
 c_D=(a_1L,a_1r+b_1),\qquad c_V=(a_2L,a_2r+b_2).
\]
After removing these fixed contents, the two affine forms in \(t\) have
positive leading coefficients and determinant
\begin{equation}
 \frac{L\Delta}{c_Dc_V}\ne0.
 \label{eq:reduced-det}
\end{equation}
Complete multiplicativity gives
\[
 \liou(D)\liou(V)=\liou(c_Dc_V)
 \liou(D/c_D)\liou(V/c_V).
\]
Tao's Theorem~2, equation~(3), printed and PDF page~3, now gives terminal
logarithmic cancellation for the reduced fixed nonparallel forms
\cite{Tao2016LogChowla}.  Treat the finitely many nonempty classes
separately.  Replacing \(r+Lt\) by \(Lt\) changes each weighted sum by
\(O(1)\).  If \(X/\omega(X)\to\infty\), the affine endpoint change produces
an admissible clock with logarithm \(\log\omega(X)+O(1)\).  If that lower
endpoint stays bounded, apply Tao's theorem with the full-prefix clock
\(\omega'(X')=X'\) after removing finitely many \(t\); then
\(\log X'\sim\log\omega(X)\).  Every subsequence has a further subsequence
of one of these two types.  Hence the fixed-\(P\) approximant is
\(o(\log\omega(X))\).

It remains to remove \(P\).  For either fixed affine form \(F\), a Boolean
union bound and the substitution \(F(n)=p^2m\) give
\begin{equation}
 \sum_{X/\omega(X)<n\le X}
 \frac{|\ind_{F(n)\ \mathrm{squarefree}}-S_P(F(n))|}{n}
 \ll_F \frac{\log\omega(X)}P+1.
 \label{eq:tail}
\end{equation}
Indeed, beyond finitely many \(n\), one has
\(n^{-1}\ll_F(p^2m)^{-1}\); the relevant harmonic \(m\)-interval has
length in logarithmic scale at most \(\log\omega(X)+O_F(1)\), and
\(\sum_{p>P}p^{-2}\ll P^{-1}\).  The two-coordinate inequality
\(|uv-UV|\le|u-U|+|v-V|\) for Boolean values makes this estimate
determinant-free.  Divide by \(\log\omega(X)\), first let \(X\to\infty\)
with \(P\) fixed, and then let \(P\to\infty\).  This proves
\eqref{eq:full-mu}.
\end{proof}

The Boolean-cutoff and content-extraction architecture follows the
TPC-137 template, which treated a frozen determinant-two setting
\cite{TPC137}.  The argument above proves the arbitrary-nonzero-determinant
form needed here, without claiming historical priority for that level of
generality.  RH-389 used TPC-137 only at determinant two
\cite{RH389}; it is not being cited as an all-\(h\) theorem.

\section{One-site and two-site local densities}
\label{sec:density}

For each prime \(p\), let
\[
 A_p^{(h)}=\{0,h\pmod {p^2}\},\qquad
 \nu_p^{(h)}=|A_p^{(h)}|,\qquad
 \tau_{p,r}^{(h)}=|\{a\in A_p^{(h)}:a\equiv r\pmod p\}|.
\]
The braces denote a set of distinct residues modulo \(p^2\).  Thus a
collision modulo \(p^2\) is counted once, whereas two distinct residues
that collide only modulo \(p\) both contribute to \(\tau_{p,r}^{(h)}\).

\begin{proposition}[Exact phase densities]
\label{prop:density}
For fixed \(q,h\ge1\),
\begin{align}
 \dd_{q,r}={}&\frac1q
 \prod_{p\nmid q}(1-p^{-2})
 \prod_{p\parallel q}\left(1-\frac{\ind_{p\mid r}}p\right)
 \prod_{p^2\mid q}\ind_{p^2\nmid r},
 \label{eq:delta-local}\\
 \thh^{(h)}_{q,r}={}&\frac1q
 \prod_{p\nmid q}\left(1-\frac{\nu_p^{(h)}}{p^2}\right)
 \prod_{p\parallel q}\left(1-\frac{\tau_{p,r}^{(h)}}p\right)
 \prod_{p^2\mid q}\ind_{r\bmod p^2\notin A_p^{(h)}}.
 \label{eq:theta-local}
\end{align}
Moreover,
\begin{gather}
 0\le\thh^{(h)}_{q,r}\le\dd_{q,r},\qquad
 \thh^{(h)}_{q,r}\le\dd_{q,r-h},
 \label{eq:cone}\\
 \sum_{r\bmod q}\dd_{q,r}=\frac6{\pi^2},\qquad
 \sum_{r\bmod q}\thh^{(h)}_{q,r}=\kap_h.
 \label{eq:totals}
\end{gather}
\end{proposition}

\begin{proof}
Impose squarefreeness only at primes \(p\le R\), with \(R\) containing
the prime divisors of \(q\).  The Chinese remainder theorem gives an exact
finite density.  If \(p\nmid q\), it removes \(\nu_p^{(h)}\) residues
modulo \(p^2\).  If \(p\parallel q\), the fixed residue modulo \(p\) has
\(p\) lifts, of which \(\tau_{p,r}^{(h)}\) are forbidden.  If
\(p^2\mid q\), the phase already fixes the residue modulo \(p^2\), giving
the last indicator.  This proves the finite truncation of
\eqref{eq:theta-local}; the one-site computation is identical with the
single forbidden residue \(0\).

The omitted-prime error is bounded by the density of integers for which
\(p^2\mid n\) or \(p^2\mid n-h\) for some \(p>R\).  The union bound is
\(O(\sum_{p>R}p^{-2})+o_N(1)\).  First \(N\to\infty\), then
\(R\to\infty\), proving both infinite products.  This is a direct
finite-CRT-plus-tail proof; Mirsky's squarefree-pattern work is historical
context, not a black-box source for the phase formula \cite{Mirsky1948}.

The cone follows by set inclusion of squarefree pairs into each one-site
event.  Summing phases removes the congruence.  At \(p\), the two forbidden
residues coincide precisely when \(p^2\mid h\); hence the total pair
density is the product \(\kap_h\) in \eqref{eq:capacity}.
\end{proof}

The collision convention is visible in small cases.  For
\((h,q,r)=(2,2,0)\), the residues \(0,2\pmod4\) are distinct but both are
\(0\pmod2\), so \(\tau_{2,0}^{(2)}=2\) and the \(p\parallel q\) factor is
zero.  For \((6,3,0)\), the analogous factor is \(1-2/3=1/3\).  In
contrast, \((h,q,r)=(4,2,0)\) gives the singleton
\(A_2^{(4)}=\{0\}\) and factor \(1/2\), while \((9,3,0)\) gives factor
\(2/3\).  These examples rule out
both loss of a modulo-\(p\) collision and double counting of a genuine
modulo-\(p^2\) collision.

\section{Diagonalization proofs}

For any fixed periodic \(\rho\), Davenport cancellation in fixed
arithmetic progressions gives
\begin{equation}
 \sum_{n\le N}\mob(n-a)\rho(n)=o(N).
 \label{eq:one-form}
\end{equation}
It also gives the masked estimates
\begin{equation}
 \sum_{n\le N}\mob(n-h)\mob(n)^2\rho(n)=o(N),\qquad
 \sum_{n\le N}\mob(n-h)^2\mob(n)\rho(n)=o(N).
 \label{eq:masked}
\end{equation}
Indeed, insert \(\mob(m)^2=\sum_{d^2\mid m}\mob(d)\), keep
\(d\le R\), and apply \eqref{eq:one-form} on the resulting finitely many
progressions.  The absolute tail is \(O(N/R+\sqrt N)\); take \(N\) and
then \(R\) to infinity \cite{Davenport1937}.  Lemma~\ref{lem:abel}
transfers \eqref{eq:one-form}, \eqref{eq:masked}, and
Proposition~\ref{prop:density} to every terminal clock.

\begin{proof}[Proof of Theorem~\ref{thm:finite-shift}]
Expand each phase polynomial.  The constant has phase density \(1/q\),
linear monomials vanish by \eqref{eq:one-form}, and diagonal squares give
\(\dd_{q,r-a_i}\).  For \(i\ne j\), apply
Theorem~\ref{thm:full-mu} to \(n-a_i,n-a_j\); their determinant is
\(a_i-a_j\ne0\).  The finite phase sum gives \eqref{eq:finite-shift}.
\end{proof}

\begin{proof}[Proof of Theorem~\ref{thm:biquadratic}]
The constant, \(x^2\), \(z^2\), and \(x^2z^2\) channels give respectively
\(1/q,\dd_{q,r-h},\dd_{q,r}\), and \(\thh^{(h)}_{q,r}\).
The \(x,z\) channels vanish by \eqref{eq:one-form}; \(xz^2,x^2z\) vanish
by \eqref{eq:masked}; and \(xz\) vanishes by
Theorem~\ref{thm:full-mu}, whose determinant is \(h\ne0\).  Summing phases
proves \eqref{eq:biquadratic}.
\end{proof}

\section{Truth-table projection and the all-\texorpdfstring{$q,h$}{q,h} charge}

For a phase table, let \(Q_r\) be the unique coordinatewise-biquadratic
interpolant of \(zf_r(x,z)\) on \(\T^2\).  Since \(Q_r(x,0)=0\), its
\(c_{00},c_{10},c_{20}\) coefficients vanish; the compiler coordinates are
\((c_{01},c_{02},c_{11},c_{12},c_{21},c_{22})\).  Theorem
\ref{thm:biquadratic} therefore gives
\begin{equation}
 L_{q,h}(f;\omega)=\sum_{r\bmod q}
 \{c_{02}(r)\dd_{q,r}+c_{22}(r)\thh^{(h)}_{q,r}\},
 \label{eq:table-limit}
\end{equation}
which is independent of \(\omega\).

Replace \(E_r\) pointwise by
\(E_r^+=E_r\cap(\T\times\{+1\})\).  At \(z=+1\) the score \(zf_r\) is
unchanged, at \(z=0\) it is zero, and at \(z=-1\) deleting a plus edge
changes the score from \(-1\) to \(+1\).  Projection therefore never
decreases a finite signed score and cannot create a composable pair.  It
maps the \(512\) tables onto the eight actions
\[
 A_r=\{x:(x,+1)\in E_r\}\subseteq\T,
\]
with \(64\) preimages per action.  Two phases \(r,r+h\) are compatible
exactly when
\begin{equation}
 A_r=\varnothing\quad\hbox{or}\quad +1\notin A_{r+h}.
 \label{eq:compatibility}
\end{equation}

Write \(\dd_r=\dd_{q,r}\) and \(\thh_r=\thh^{(h)}_{q,r}\).  Exact
interpolation gives the eight limiting weights:
\begin{center}
\small
\begin{tabular}{@{}c@{\qquad}c@{\qquad}c@{\qquad}c@{}}
\toprule
Action \(A\)&\(c_{02}\)&\(c_{22}\)&\(w_r(A)\)\\
\midrule
\(\varnothing\)&0&0&0\\
\(\{-1\}\)&0&\(1/2\)&\(\thh_r/2\)\\
\(\{0\}\)&1&\(-1\)&\(\dd_r-\thh_r\)\\
\(\{-1,0\}\)&1&\(-1/2\)&\(\dd_r-\thh_r/2\)\\
\(\{+1\}\)&0&\(1/2\)&\(\thh_r/2\)\\
\(\{-1,+1\}\)&0&1&\(\thh_r\)\\
\(\{0,+1\}\)&1&\(-1/2\)&\(\dd_r-\thh_r/2\)\\
\(\{-1,0,+1\}\)&1&0&\(\dd_r\)\\
\bottomrule
\end{tabular}
\end{center}
The \(512\)-table lag-two setting originates in RH-378; the positive-current
projection and eight-action table are inherited from RH-389
\cite{RH378,RH389}.  All-\(h\) compatibility, density, and charge are proved
here.

Use the baseline \(B_0=\{-1,0\}\) and
\(H_r=\dd_r-\thh_r/2\).  An action without \(+1\) has weight at most
\(H_r\), while an action containing \(+1\) gains at most \(\thh_r/2\)
above \(H_r\).  If \(+1\in A_r\), compatibility forces
\(A_{r-h}=\varnothing\).  The loss at that predecessor pays the gain since
\begin{equation}
 H_{r-h}-\frac{\thh_r}{2}
 =\frac{\dd_{r-h}-\thh_{r-h}}2+
  \frac{\dd_{r-h}-\thh_r}2\ge0.
 \label{eq:charge}
\end{equation}
Translation \(r\mapsto r-h\) is a permutation of \(\Zq\), so distinct
plus phases have distinct charged predecessors.  It has \((q,h)\) cycles,
each of length \(q/(q,h)\).  If \(q\mid h\), every cycle is a self-loop
and \eqref{eq:compatibility} forces the plus-phase set empty.  For example,
\((q,h)=(6,4)\) has two cycles of length three; no coprimality of \(q,h\)
is needed.  Summing \eqref{eq:charge} and using \eqref{eq:totals} yields
\begin{equation}
 \sum_{r\bmod q}w_r(A_r)\le
 \sum_{r\bmod q}H_r=\frac6{\pi^2}-\frac{\kap_h}{2}.
 \label{eq:signed-bound}
\end{equation}
The constant baseline attains equality.

Finally reflect inputs by \(f^\rho(x,z)=f(-x,-z)\).  Safety is preserved,
and the six compiler coefficients transform with signs
\((+,-,-,+,+,-)\).  The surviving \(c_{02},c_{22}\) channels therefore
change sign.  Baseline table \(36\) attains the positive value in
\eqref{eq:signed-bound}; its reflection is table \(72\) and attains the
negative value.  This proves the absolute maximum in
Theorem~\ref{thm:capacity}.

\section{Square-divisor landscape}

If \(h\) is squarefree, no local factor in \(\kap_h\) is replaced, so
\(\kap_h=\kap_\star\).  If \(p^2\mid h\), the factor
\(1-2p^{-2}\) is replaced by the strictly larger \(1-p^{-2}\).  Hence
\(\kap_h=\kap_\star\) exactly for squarefree \(h\), proving the maximum
claim in \eqref{eq:range}.

For every finite \(h\), at least one prime-square factor remains strict,
and comparison with \(\prod_p(1-p^{-2})=6/\pi^2\) gives
\(\kap_h<6/\pi^2\).  Thus \(\Glog(q,h)>3/\pi^2\).  On the other hand,
for
\[
 h_y=\left(\prod_{p\le y}p\right)^2,
\]
the factors through \(y\) equal those of \(6/\pi^2\), while the tail ratio
tends to one.  Consequently \(\kap_{h_y}\to6/\pi^2\) and
\(\Glog(q,h_y)\to3/\pi^2\).  The infimum is therefore not attained.

\section{Executable artifact and source boundary}

The Stage-1 certificate contains \(640\) exact rows,
\[
 640=512+8+64+8+9+7+12+8+6+6,
\]
covering truth tables, actions, compatibility, charge, monomials,
determinant completion, local densities, the lag spectrum, finite shifts,
and firewalls.  It exhausts all \(512^2\) ordered table pairs, finding
\(3375\) compatible pairs and zero projection, reflection, involution,
parity, or interpolation failures.  It also rejects \(24\) semantic
mutations under a builder-independent verifier.  The canonical certificate
has \(220832\) bytes and SHA-256
\begin{center}\small\ttfamily
614297795d4d4dfeadfb5667d3e0d405\allowbreak
d04fbe8e07e9d87a743faed9cb267a96.
\end{center}
These checks reproduce finite algebra; they do not prove
Theorem~\ref{thm:full-mu} or Proposition~\ref{prop:density}.

The immutable closure is rebound to RH-389 release commit
\begin{center}\small\ttfamily
8b1a875b4bbefd955a419593951ce2d09987ac6f.
\end{center}
Its Git groups have sizes \(95,8,3\), with ordered group digests
\begin{center}\small\ttfamily
8a674e5d60237b4463e1f68ef79965633ed11a4d098957e5be44e05f471174cb,\\
87bbdb455fa5217404d863d1054b4ae408de69da6701ff42ac439ec9bfe1605c,\\
a03e4ba7d8b5b054acc95288c70e753bf22b82bb1957f635891b28682c67840e.
\end{center}
The all-\(106\) Git digest is
\begin{center}\small\ttfamily
3b32865a14618a605915beb8eab6432b048fca49718b69519697ef861cbe650f,
\end{center}
and adding three remote logical locks gives \(109\) objects with digest
\begin{center}\small\ttfamily
39bf8e9030b511e85fdf26a7c71722c3e4be5bc74bc738aa253bbc29c94517f9.
\end{center}

Among the three remote locks, only Tao is an analytic proof input.  Its
official Cambridge version of record locates Theorem~2, equation~(3), on
printed/PDF page~3 and is licensed CC BY~4.0; project policy nevertheless
does not vendor the PDF.  The inherited Johnston--Yang and Maynard locks
are closure-only and are not used in the proof.  No article-specific
redistribution grant is established for either of those two payloads, so
both remain remote and unvendored.  All three offline verifiers default to
zero network requests.  The publication tree excludes all five external
payload identities (the Johnston--Yang PDF, source archive, and extracted
source, the Maynard PDF, and the Tao PDF).  No new remote source was added.

Relative to the RH-378/RH-389/TPC-137 templates, the two local proof
completions are the arbitrary-nonzero-determinant Boolean cutoff argument
and the all-\(h\) phasewise CRT formula.  This is a package-level provenance
statement, not a global priority claim.

\section{Limitations and declarations}

The quantifiers are fixed-data quantifiers.  There is no uniform theorem
for \(q=q(X)\), \(h=h(X)\), a growing shift family, or an effective rate.
The paper proves terminal logarithmic averages, not ordinary Ces\`aro
averages.  The finite-shift theorem stops at total degree two.  The
coordinatewise-biquadratic compiler concerns one fixed lag and does not
assert degree-three multi-coordinate or interacting-multiple-lag truth-table
laws.  The maximum in \eqref{eq:capacity-def} is taken only after each
fixed-table limit; no maximum-before-limit or adaptive selector is present.
Nothing here constructs an operator, trace formula, zero model, or proof of
the Riemann Hypothesis, and all five physical gates remain open.

\paragraph{Data availability.}
No new empirical data were created or analyzed.  The exact certificate,
result, recursively closed JSON Schema, source-lock records, and tests are
included in the paper package.  External PDFs and source archives are not
vendored.

\paragraph{Ethics declaration.}
This theoretical and computational study involved no human participants,
animals, personal data, or clinical intervention.  Institutional ethics
approval was not applicable.

\paragraph{Author contributions.}
RH research program: Conceptualization, Methodology, Formal analysis,
Software, Validation, Writing--original draft, and Writing--review and
editing.

\paragraph{Conflict of interest.}
The author declares no conflict of interest.

\paragraph{Funding.}
No external funding was received for this work.

\paragraph{AI-use disclosure.}
AI-assisted research tools were used for symbolic enumeration, proof and
source-audit orchestration, manuscript drafting, and formatting.  The
theorem statements, computations, citations, and claim boundaries were
checked against the frozen sources and executable artifacts.  The author
takes responsibility for the accuracy and integrity of the work.

\begingroup
\small
\bibliographystyle{plainnat}
\bibliography{references}
\endgroup

\end{document}
